\subsection{Automated LTL Verification of Buildable System Modes}
\label{sec:case-solver}

Once the finite-capacity technology catalog $TYR_N$ is fixed, the boundary specification $\Phi_{IOR_N}$ of Equation~\eqref{eq:rle-fin-recurrence} compiles into six families of propositional LTL clauses. Writing $f\in\{0,1\}$ for the boundary input bit $f(t)$, $y=(y_{\mathrm{bit}},y_{\mathrm{cnt}})\in\{0,1\}\times\{0,\dots,N\}$ for the boundary register output pair $y(t)$ (using $y=(b,k)$ as shorthand for $y_{\mathrm{bit}}=b \wedge y_{\mathrm{cnt}}=k$), and $c\in\{0,1\}$ for the sense/act head phase $c(t)$, Table~\ref{tab:rle-ltl-clauses} summarizes the compiled clause families comprising $\Phi_{IOR_N}$ along with their count and derivation from the boundary recurrence~\eqref{eq:rle-fin-recurrence}. Here $X^k$ denotes $k$ consecutive next-step operators ($X^3 \equiv XXX$), reflecting the three-tick end-to-end pipeline latency. Together, these six families contribute $1+1+2N+2+2N+2N = 6N+4$ propositional LTL clauses at capacity $N$.

\begin{table*}[t]
\caption{Propositional LTL compilation of the finite-tape boundary specification $\Phi_{IOR_N}$ from recurrence~\eqref{eq:rle-fin-recurrence}.}
\label{tab:rle-ltl-clauses}
\renewcommand{\arraystretch}{1.25}
\footnotesize
\begin{tabularx}{\textwidth}{@{}l >{$}l<{$} c X@{}}
\toprule
\textbf{Clause} & \textbf{Propositional LTL Formula} & \textbf{Count} & \textbf{Derivation from Boundary Requirement $IOR_N$} \\
\midrule
$\varphi_{\mathrm{init}}$ & (y{=}(0,0)) \wedge X(y{=}(0,0)) \wedge X^2(y{=}(0,0)) & 1 & Forces the idle pair $(0,0)$ across the three-tick pipeline startup window ($t=0,1,2$) before the first encoded symbol emerges. \\
$\varphi_{\mathrm{mode}}$ & G\bigl((c{=}0 \leftrightarrow X(c{=}1)) \wedge (c{=}1 \leftrightarrow X(c{=}0))\bigr) & 1 & Enforces the two-phase alternating clock discipline (sense: $c=0$, act: $c=1$) governing the synchronous Moore pipeline. \\
$\varphi_{\mathrm{hold}}$ & G\bigl((c{=}1 \wedge y{=}(b,k)) \rightarrow X(y{=}(b,k))\bigr),\ b\in\{0,1\}, k\in\{1,\dots,N\} & 2N & Requires the output register to hold a valid run pair $(b,k)$ across the second half of each two-phase cycle, matching the step-synchronous spec. \\
$\varphi_{\mathrm{start}}$ & G\bigl((c{=}0 \wedge f{=}b \wedge X(y_{\mathrm{cnt}}{=}0)) \rightarrow X^3(y{=}(b,1))\bigr),\ b\in\{0,1\} & 2 & Translates the base case $\ell_N(f,0)=1$ of~\eqref{eq:rle-fin-recurrence}, emitting the initial run pair $(b,1)$ after the three-tick pipeline latency ($X^3$). \\
$\varphi_{\mathrm{inc}}$ & G\bigl((c{=}0 \wedge f{=}b \wedge X(y{=}(b,k))) \rightarrow X^3(y{=}(b,\min(k{+}1,N)))\bigr),\ k\le N & 2N & Translates the matching-symbol recurrence ($\ell_N(f,k{+}1)=\min(\ell_N(f,k){+}1,N)$ for $f(2k{+}2)=f(2k)$), incrementing the count with capacity saturation. \\
$\varphi_{\mathrm{reset}}$ & G\bigl((c{=}0 \wedge f{=}b \wedge X(y{=}(\neg b,k))) \rightarrow X^3(y{=}(b,1))\bigr),\ k\le N & 2N & Translates the run-breaking recurrence ($\ell_N(f,k{+}1)=1$ for $f(2k{+}2)\neq f(2k)$), resetting the counter to $1$ and updating the active symbol bit. \\
\midrule
\multicolumn{2}{@{}l}{\textbf{Total Propositional Clauses ($N$-capacity):}} & \textbf{6N+4} & 16 ($N{=}2$), 22 ($N{=}3$), 28 ($N{=}4$) across all six families. \\
\bottomrule
\end{tabularx}
\end{table*}

We implemented an automated verification workflow in Python (3.12) on Ubuntu 24.04 under Windows Subsystem for Linux 2 (WSL2). Given a candidate component catalog $TYR_N$ and coupling recipe $SCR_N$, the workflow constructs the coupled resultant $Z_{b,N}=\mathrm{RSY}(SCR_N)$, isolates the reachable operational system mode $Z_{m,N}\in\mathrm{SYSMO}(Z_{b,N})$ from the initial reset state, and checks the boundary LTL specification $\Phi_{IOR_N}$ (Table~\ref{tab:rle-ltl-clauses}) by explicit-state exploration of that mode. Table~\ref{tab:solver-results} reports those $\Phi_{IOR_N}$ proof-of-concept results. When a prescribed reference $Z_{\mathrm{ref},N}$ is in view, Theorem~\ref{thm:dynamics-characterization} still applies: $Z_{m,N}\models_{\mathrm{dyn}}\Phi_{\mathrm{dyn}}(Z_{\mathrm{ref},N})$ is the existential SAT query $\exists(h_S,h_I,h_O)$ (functionality, readout/step preservation, and surjectivity), which can be discharged by a propositional encoding in Z3~\cite{deMouraZ32008} that discovers the maps rather than requiring a hand-written architectural projection. In the nominal four-component architecture, the input tape ($|S_0|=2$), sense/act head ($|S_1|=2$), output register ($|S_2|=2(N{+}1)$), and instruction table ($|S_3|=4(N{+}1)$) form a Cartesian product resultant state space of size $|S_{\mathrm{resultant}}|=32(N{+}1)^2$, whereas synchronous coupling restricts the reachable operational mode $Z_{m,N}$ to a much smaller invariant sub-system $|S_{\mathrm{mode}}|$.

\begin{table*}[t]
\caption{Automated boundary LTL verification results for finite-tape buildable RLE resultants $Z_{b,N}=\mathrm{RSY}(SCR_N)$ and their operational system modes $Z_{m,N}\in\mathrm{SYSMO}(Z_{b,N})$ against $\Phi_{IOR_N}$.}
\label{tab:solver-results}
\renewcommand{\arraystretch}{1.25}
\footnotesize
\begin{tabularx}{\textwidth}{@{}X r r r r c c r@{}}
\toprule
\textbf{Candidate Design} & \textbf{Cap $N$} & \textbf{$|S_{\mathrm{resultant}}|$} & \textbf{$|S_{\mathrm{mode}}|$} & \textbf{Clauses} & \textbf{Verdict} & \textbf{Violated} & \textbf{Wall time (ms)} \\
\midrule
\csname @@input\endcsname generated/solver-results.tex
\bottomrule
\end{tabularx}
\end{table*}

The nine candidate designs in Table~\ref{tab:solver-results} are organized into three evaluation groups that test capacity scaling, subsystem refactoring invariance, and fault and timing-mismatch diagnosis.

The first three rows (\emph{Baseline coupled RLE} for $N=2,3,4$) evaluate how verification scales as the hardware counter capacity $N$ of the technology catalog $TYR_N$ increases. Each design is constructed by interconnecting the nominal input tape ($Z_0$), sense/act head ($Z_1$), output register ($Z_2$), and instruction table ($Z_3$) through the coupling recipe $SCR_N$ of Figure~\ref{fig:rle-architecture}, with the count registers in $Z_2$ and $Z_3$ bounded to $\{0,\dots,N\}$. Increasing $N$ from $2$ to $4$ expands the coupled resultant state space $|S_{\mathrm{resultant}}|$ from $288$ to $800$ states and the reachable operational mode $|S_{\mathrm{mode}}|$ from $31$ to $55$ states ($12N+7$). All three designs pass all $6N+4$ LTL clauses ($16$, $22$, and $28$ clauses, respectively). By Theorem~\ref{thm:mode-ior-bridge}, each PASS verdict on $\Phi_{IOR_N}$ certifies that $(D_{m,N},Z_{b,N})\in\mathrm{CTL}(IOR_N,TYR_N)$ for the mode-as-functional-design $D_{m,N}=(Z_{m,N},s_{b,0},T)$, without requiring manual construction of a state-space map $h_S$.

The next two rows (\emph{Phase-inverted instruction table} and \emph{Parity-instrumented output register}, both at $N=3$) test whether boundary and dynamics-fragment verification remain invariant when subsystem teams refactor internal state representations or elaborate components with auxiliary diagnostic state. Both candidates modify a single component within the nominal $N=3$ coupling recipe $SCR_3$. The \emph{phase-inverted instruction table} design replaces $Z_3$ with an isomorphic component $Z_3'$ whose internal sense/act phase bit is inverted ($0\leftrightarrow 1$) together with its transition and readout tables, leaving $|S_{\mathrm{resultant}}|=512$ and $|S_{\mathrm{mode}}|=43$ unchanged. The \emph{parity-instrumented output register} design replaces $Z_2$ with an elaborated register $Z_2^+$ whose state space $S_2\times\{0,1\}$ appends an internal diagnostic parity bit that toggles on each processed input bit while emitting the same boundary pair $(\text{bit},\text{count})$; this doubles the resultant state space from $512$ to $1024$ states and expands the reachable operational mode from $43$ to $79$ states. Both designs pass all clauses of $\Phi_{IOR_3}$ without any change to the requirement specification. Whereas classical T3SD verification would in general require rewriting an explicit state homomorphism $h_S$ to decode the inverted phase bit in $Z_3'$ or project away the diagnostic parity state in $Z_2^+$, Theorems~\ref{thm:encoding-invariance} and~\ref{thm:elaboration-coupling} ensure that boundary temporal satisfaction absorbs internal re-encodings and state elaborations automatically, certifying that both refactored architectures remain in $\mathrm{CTL}(IOR_3,TYR_3)$.

The final four rows evaluate diagnostic rejection when a component in $TYR_3$ violates its behavioral contract or when the specification clock is misaligned with the hardware pipeline. The \emph{stuck sense/act head} design replaces $Z_1$ with a faulty head whose transition map fails to toggle and remains fixed in its initial phase; the checker rejects the resultant on the phase-alternation clause $\varphi_{\mathrm{mode}}$. The \emph{frozen input tape} design replaces $Z_0$ with a faulty tape whose transition map ignores incoming external bits and permanently retains its initial $0$ bit; because a $1$ bit is never latched, the reachable operational mode collapses from $43$ states to $9$ states and the checker rejects the design on the initial symbol-acquisition clause $\varphi_{\mathrm{start}}$. The \emph{non-resetting instruction table} design replaces $Z_3$ with a faulty table that continues incrementing the run count (up to $N=3$) when the input bit changes instead of resetting the count to $1$; although the resultant reaches $43$ mode states and satisfies the startup, hold, and same-bit increment clauses, the checker rejects it on the run-reset clause $\varphi_{\mathrm{reset}}$. Finally, the \emph{unshifted single-tick specification} candidate pairs the nominal, fault-free $N=3$ coupled resultant ($|S_{\mathrm{resultant}}|=512$, $|S_{\mathrm{mode}}|=43$) against an idealized single-tick requirement that expects the boundary output to update on every clock tick rather than holding across the two-phase sense/act cycle. The checker rejects this pairing on $\varphi_{\mathrm{hold}}$, confirming empirically the timing-alignment requirement discussed in Section~\ref{sec:case-machine}: even when the underlying hardware components are sound, no lock-step Wymore homomorphism exists unless the specification and the buildable resultant share a matched clock granularity.
